\documentclass[10pt,a4paper]{amsart}
\usepackage[margin=19mm]{geometry}
\usepackage{amsmath,amssymb,mathtools}
\usepackage[scr=rsfs,cal=euler]{mathalfa}
\usepackage{xfrac}
\usepackage[unicode,hidelinks,bookmarks=false]{hyperref}
\newcommand{\ie}{i.e.\ }
\newcommand{\cf}{cf.\ }
\newcommand{\resp}{respectively\ }
\newcommand{\Subsection}{\subsection}
\providecommand{\nobreakdash}{\nobreak\hbox{-}}
\setlength{\emergencystretch}{2em}
\multlinegap0pt
\usepackage{amssymb}
\usepackage{mathtools}
\usepackage{xcolor}
\newtheorem{theorem}{Theorem}
\newcommand{\Gr}{\operatorname{Gr}}
\newcommand{\Span}{\operatorname{span}}
\newcommand{\chg}[1]{#1}
\title{Determinantal Kernel Schemes of Matrix Nets and Applications to Positive Maps}
\author{Trung Hoa Dinh, Minh Toan Ho, Cong Trinh Le, Trung Dung Vuong}
\date{Source: arXiv:2609.09675v1 (CC BY 4.0)}
\begin{document}
\maketitle

\noindent\emph{Source and licence.} Determinantal Kernel Schemes of Matrix Nets and Applications to Positive Maps. Version arXiv:2609.09675v1 (CC BY 4.0), distributed by arXiv under the Creative Commons Attribution 4.0 International licence, http://creativecommons.org/licenses/by/4.0/, as recorded in the arXiv licence metadata for this exact version and shown on the abstract page https://arxiv.org/abs/2609.09675v1. Full licence notice: \texttt{licenses/arxiv-2609.09675-CC-BY-4.0.txt}.\par\smallskip

\noindent\textit{Editorial scope.} Counted unit: the article's theorem thm:global-components (source0.tex lines 317--319). The article's complete original proof of it is delivered with it (source0.tex lines 321--377, one \texttt{proof} environment of 57 lines). No mathematics is written, repaired or re-derived: the statement and the proof are the article's own lines, in the article's own notation, and no retained line is substituted or re-flowed.  Retained uncounted as the source-local context the proof needs: lines 104--111; lines 115--125; lines 297--299.  Nothing was omitted from any retained span, so the package carries no omission record.  No retained line contains a citation, so the wrapper carries no bibliography and no bibliography entry.  No retained line contains a figure environment, an image-inclusion command or drawing code, so no image is delivered and the package declares no asset dependency.  Editorial formatting only, in addition: an AMS \texttt{amsart} preamble that restates the article's own declarations and macros used by the retained lines, namely the environments \texttt{theorem} on separate counters, together with the standard packages those lines need.  The retained environments print in this document as Theorem 1 in order of appearance, whereas the article prints the same items with the numbering of its own full document.  The complete native source of the article as distributed in the arXiv e-print for this exact version is bundled unchanged as \texttt{sources/209843/source0.tex}.
\begin{theorem}\label{thm:main-positive-dimensional}
Let $a,b\ge2$, let $K\subset M_{a,b}$ have dimension three, and put
\[
Z_K=\mathbb P(K)\cap\Sigma_{a,b},\qquad
\Sigma_{a,b}=\mathbb P^{a-1}\times\mathbb P^{b-1}\subset\mathbb P(M_{a,b}).
\]
If $\dim Z_K\ge1$, then $Z_K$ is exactly one of the following scheme types: a ruling plane, a smooth conic, two reduced Segre lines from opposite rulings, a reduced Segre line, a reduced Segre line with one reduced point off the line, or a Segre line with {one embedded residual point contributing length one}. A doubled Segre line cannot occur outside a ruling linear space. In the \chg{smooth conic} case, $K$ lies in a $2\times2$ compression and is equivalent to the traceless hyperplane of $M_2(\mathbb C)$.
\end{theorem}

\begin{theorem}\label{thm:components}
The locus
\[
\mathcal P_{a,b}=\{K\in\Gr(3,M_{a,b}):\dim(\mathbb P(K)\cap\Sigma_{a,b})\ge1\}
\]
has exactly the three irreducible components $\mathcal L_A$, $\mathcal L_B$, and $\mathcal C$, with dimensions
\[
ab+a+2b-8,\qquad ab+2a+b-8,\qquad 2a+2b-5.
\]
Their pairwise intersections coincide with an irreducible locus of dimension $2a+2b-6$. The compression component is smooth, and this common intersection is its relative Segre divisor.
\end{theorem}

\begin{theorem}\label{thm:positive-plane-sections}
Let $K\subset M_{a,b}$ have dimension three. If $S=\mathbb P(K)$ meets $\Sigma_{a,b}$ in positive dimension, then the underlying set is one of the four alternatives in Theorem~\ref{thm:main-positive-dimensional}: a ruling plane, a smooth conic, two \chg{Segre lines from opposite rulings}, or one Segre line with at most one residual point. The complete scheme structures are precisely the six types in Theorem~\ref{thm:main-positive-dimensional}.
\end{theorem}

\begin{theorem}\label{thm:global-components}
The conclusions and dimension formulas of Theorem~\ref{thm:components} hold.
\end{theorem}

\begin{proof}
The classification in Theorem~\ref{thm:positive-plane-sections} gives
\[
\mathcal P_{a,b}=\mathcal L_A\cup\mathcal L_B\cup\mathcal C.
\]
Lines from the first ruling are parametrized by
\[
F_A=\mathbb P(A)\times\Gr(2,B),
\qquad
\dim F_A=(a-1)+2(b-2)=a+2b-5.
\]
For a fixed line with underlying \chg{two dimensional} vector space $L$, the \chg{three dimensional} subspaces $K\supset L$ are parametrized by
\[
\Gr(1,(A\otimes B)/L)\simeq\mathbb P^{ab-3}.
\]
Thus the incidence variety is projective and irreducible of dimension $ab+a+2b-8$, and its image $\mathcal L_A$ is closed and irreducible. Its projection is generically finite: the embedded $3\times3$ example
\[
K=\Span\{E_{11},E_{12},E_{22}+E_{33}\}
\]
viewed inside $M_{a,b}$ has exactly the displayed Segre line from the first ruling. Hence
\[
\dim\mathcal L_A=ab+a+2b-8.
\]
The other formula is symmetric.

For the compression component, let
\[
B_0=\Gr(2,A)\times\Gr(2,B)
\]
and let $\mathcal U,\mathcal Z$ be the tautological rank two bundles. The projective bundle
\[
X=\mathbb P((\mathcal U\otimes\mathcal Z)^*)\longrightarrow B_0
\]
parametrizes triples $(U,V,[\ell])$ and their hyperplanes $K=\ker\ell\subset U\otimes V$. The induced morphism $X\to\Gr(3,A\otimes B)$ has image $\mathcal C$.

For every such $K$, its left and right supports have dimension exactly two: if, for instance, the left support had dimension one, then $K$ would lie in a \chg{two dimensional} space $\mathbb Cu\otimes V$, impossible because $\dim K=3$. Thus the supports recover $U$ and $V$ uniquely. Indeed, the left support is the image of the contraction map $K\otimes B^*\to A$, and the right support is the image of $K\otimes A^*\to B$. Their ranks are constantly two on $\mathcal C$, so these images define regular \chg{Grassmannian valued} support maps. They give an inverse to $X\to\mathcal C$. Hence
\[
\mathcal C\simeq X,
\qquad
\dim\mathcal C=2(a-2)+2(b-2)+3=2a+2b-5,
\]
and $\mathcal C$ is smooth.

{Let $\mathcal C_{\mathrm{red}}$ denote the locus of reducible conic sections in $\mathcal C$.}
Inside each $\mathbb P^3$ fiber, singular plane sections of $\mathbb P^1\times\mathbb P^1$ are exactly the \chg{rank one} functionals, forming $\mathbb P^1\times\mathbb P^1$. Therefore
\[
{\mathcal C_{\mathrm{red}}}
=\mathbb P(\mathcal U^*)\times_{B_0}\mathbb P(\mathcal Z^*)
\]
is smooth, irreducible, and has dimension $2a+2b-6$. A compression plane contains a Segre line exactly when its quadric section is reducible, so
\[
\mathcal L_A\cap\mathcal C
=\mathcal L_B\cap\mathcal C
={\mathcal C_{\mathrm{red}}}.
\]
If a plane contains one line from each ruling, the two lines intersect and span the tangent plane of the corresponding $2\times2$ Segre quadric. Hence $\mathcal L_A\cap\mathcal L_B={\mathcal C_{\mathrm{red}}}$. Finally, none of the three irreducible loci contains another, as witnessed by a \chg{smooth conic compression plane}, the \chg{single line examples} above, and their transposes.
\end{proof}

\end{document}
